\section{Proof of the certified root enclosure}
\label{boundary:app:step-parametrix:rr}

\subsection{A full-space residual and the existence witness}
\label{v12:sec:gap-certificate}

Let $u\in L^2(-1,1)$, let $\widetilde u$ be its extension by zero, and
let $\mathcal H$ be the full-line Hilbert transform with kernel
$[\pi(x-y)]^{-1}$.  Put
\[
 q_u=\langle u,Q_0u\rangle,\qquad
 c_u=\operatorname{Re}\langle Q_0u,Ju\rangle.
\]
Choose any finite orthonormal family $e_j$ supported outside $(-1,1)$
and set $E_u=\sum_j|\langle e_j,\mathcal H\widetilde u\rangle|^2$.
For real $a$, Hilbert-transform isometry and Bessel's inequality give
\begin{equation}\label{v12:eq:outside-residual-certificate}
 \|(Q_0-aJ)u\|^2
 \le q_u-\tfrac14E_u-2ac_u+a^2\|u\|^2=:U_u(a).
\end{equation}
Indeed $Q_0=(I+iH)/2$, $H^*=-H$, and
$\|\mathcal H\widetilde u\|^2=\|u\|^2$, so
\[
 \|Q_0u\|^2=q_u-\tfrac14
       \|\mathbf1_{\mathbb R\setminus(-1,1)}\mathcal H\widetilde u\|^2.
\]
For $q_u>0$ let $f=V_0u$, where $V_0$ is the restriction of inverse
Fourier transformation to $(0,\infty)$.  Then $\|f\|^2=q_u$ and
\[
 (B_0-a)f=V_0J(Q_0-aJ)u,
 \qquad
 \operatorname{dist}(a,\operatorname{spec}B_0)
       \le\sqrt{U_u(a)/q_u}.
\]
Here the norm of $V_0J$ is at most one.  Thus an interval containing
$[a-\sqrt{U_u(a)/q_u},a+\sqrt{U_u(a)/q_u}]$ and lying in an open
central gap contains a discrete physical eigenvalue.

Take $u$ constant on cells in a $J$-invariant partition and the $e_j$
to be normalized indicators of disjoint exterior cells. The Hilbert
cell formula \eqref{rr:eq:double-cell-logarithms} evaluates $q_u,c_u,E_u$
exactly, including principal values. Omitting the other exterior modes
preserves the upper-bound direction; no tail quadrature is required.

\begin{proof}[Existence in Theorem~\ref{boundary:thm:negative-exclusion:rr}]
The exact rational trial data in \texttt{anc/boundary\_gap\_witness.json}
specify a partition of $(0,1)$, its exact translate by $-1$, complex
coefficients on the resulting 512 normalized indicator functions,
and disjoint exterior indicator functions. Write $u_0$ for this trial.
Its common positive parameter is
\begin{equation}\label{rr:eq:positive-trial}
 a_0=\frac{3750237663924847}{9007199254740992}.
\end{equation}
The independent checker \texttt{anc/verify\_boundary\_gap\_certificate.py}
evaluates the displayed double-cell logarithmic integrals with
outward-rounded Arb ball arithmetic \cite{Johansson2017}.
At both 192 and 256 bits it verifies
the strict rational bounds
\[
 0.41635<a_0<0.41637,\qquad q_{u_0}>0.32783,\qquad
 0<U_{u_0}(a_0)<0.000038641,
\]
and directly verifies $\sqrt{U_{u_0}(a_0)/q_{u_0}}<0.01086$.
All inputs are exact rationals and all rounding errors are enclosed;
no eigenvalue solver or quadrature estimate enters this check.
Thus the residual interval lies strictly inside $(0.405,0.428)$,
which is contained in $(0,g)$.  Equation
\eqref{v12:eq:outside-residual-certificate} and the spectral theorem
force a spectral point there.  The reference gap makes it a discrete
eigenvalue, and Theorem~\ref{boundary:thm:one-cell-flow} makes it simple.
\end{proof}


\subsection{The negative-window model}

We prove the exclusion at a single frequency parameter and then use a
resolvent bound.  Put
\[
 a=-\frac{833}{2000},\qquad P_a=Q_0-aJ,
 \qquad c_0=\frac\pi2,
 \qquad
 \chi(t)=\begin{cases}\cos(c_0t),&0\le t\le1,\\0,&t>1.\end{cases}
\]
Identify $L^2(-1,1)$ with $L^2((0,1);\mathbb C^2)$ by
$f_1(t)=u(t)$, $f_2(t)=u(t-1)$.  On four copies of $L^2(0,\infty)$
retain the endpoint model $Q_e,J_e,M_a=Q_e-aJ_e$ of
Section~\ref{v12:app:mellin-proofs}, at phase zero.  Define
\[
 Ef(t)=\chi(t)(f_1(t),f_2(t),f_1(1-t),f_2(1-t))^T
\]
with zero extension outside $(0,1)$.  Since
$\chi(t)^2+\chi(1-t)^2=1$, one has $E^*E=I$.  In particular no
interior model is needed for this partition.

\subsection{The exact model and the residual criterion}

The intertwining error is $\Delta=EQ_0-Q_eE$.  For later use set
\[
 D_0(t,s)=\frac{\chi(t)-\chi(s)}{t-s},\qquad
 D_1(t,s)=\frac{\chi(t)-\chi(s)}{t+s},\qquad
 D_2(t,s)=\begin{cases}\dfrac{\chi(t)}{1-t+s},&t<1,\\0,&t\ge1.
 \end{cases}
\]
For $t>0$, $0<s<1$, its complete kernel is
\begin{equation}\label{rr:eq:cosine-defect}
 \Delta(t,s)=\frac{i}{2\pi}
 \begin{pmatrix}
 D_0(t,s)&D_1(t,1-s)\\
 -D_2(t,s)&D_0(t,s)\\
 -D_0(t,1-s)&D_2(t,1-s)\\
 -D_1(t,s)&-D_0(t,1-s)
 \end{pmatrix}.
\end{equation}
Divided differences take their continuous diagonal values.  The unmatched
corner quotients are bounded because $\chi(t)=\sin(c_0(1-t))$ near one.
The entries are bounded on bounded rectangles and are $O(t^{-1})$ at
infinity.  Thus $\Delta\in\mathfrak S_2$.  The Lipschitz cutoff suffices
for this assertion and the intertwining identity; no differentiated
endpoint expansion is used here.

At the single parameter $a$ the model satisfies
\begin{equation}\label{rr:eq:model-inverse-bound}
 \|M_a^{-1}\|\le5.
\end{equation}
Here is an elementary verification of the constant.  Write $p=a^2$,
$d=qr\in[0,1/4]$, and let $m$ be the matrix in
\eqref{v12:eq:endpoint-symbol}.  For $w=z(z-1)$, direct expansion gives
\[
 \det(m-zI)=w^2-2pw+p^2-p+d(w+2p).
\]
Since $4/25<p<19/100$, this determinant at $z=1/5$ is at most
$p^2-9p/50-9/625\le-71/10000$, and at $z=-1/5$ it is at most
$p^2-49p/50+147/1250\le-31/10000$.
At $q=0$ the eigenvalues are $-|a|,|a|,1-|a|,1+|a|$.
Exactly one is below each of $-1/5$ and $1/5$.
Continuity as $q$ ranges over $[0,1]$ shows that no eigenvalue enters
$[-1/5,1/5]$, proving \eqref{rr:eq:model-inverse-bound}.

Let $C_\infty$ denote the Carleman operator with kernel
$[\pi(t+s)]^{-1}$.  The $(1,4)$ block of $Q_e$ is an orthogonal
projection: its multiplier has diagonal entries $q,r$, off-diagonal
entries $ic,-ic$, and $q+r=1$, $c^2=qr$.
The two singleton channels satisfy $q^2=q-c^2$, $r^2=r-c^2$.
Since $2c$ is the multiplier of $C_\infty$, one obtains the exact identity
\begin{equation}\label{rr:eq:model-square}
 Q_e^2=Q_e-\tfrac14\operatorname{diag}
                    (0,C_\infty^2,C_\infty^2,0).
\end{equation}

Choose any finite step-function projection $\Pi$ on $L^2(0,\infty)$,
and any source step-function projection $\Pi_s$ on $L^2(0,1)$.
Let $W,D:L^2((0,1);\mathbb C^2)\to
L^2((0,\infty);\mathbb C^4)$ vanish on
$(\Pi_s\otimes I_2)^\perp$ and have their outputs in
$\Pi\otimes\mathbb C^4$.  Put
$Q_\Pi=(\Pi\otimes I_4)Q_e(\Pi\otimes I_4)$ and
$C_\Pi=\Pi C_\infty\Pi$.
All inner products below are Hilbert--Schmidt inner products.  Define
\begin{align}
 \mathcal R(W,D)^2={}&
 \langle W,Q_\Pi W\rangle
 -\tfrac14\bigl(\|C_\Pi W_2\|_2^2+\|C_\Pi W_3\|_2^2\bigr)
 -2a\operatorname{Re}\langle Q_\Pi W,J_eW\rangle
 +a^2\|W\|_2^2\nonumber\\
 &-2\operatorname{Re}\langle(Q_\Pi-aJ_e)W,D\rangle+\|D\|_2^2.
 \label{rr:eq:residual-quadratic}
\end{align}
Then
\begin{equation}\label{rr:eq:full-model-residual}
 \|M_aW-D\|_2^2\le\mathcal R(W,D)^2.
\end{equation}
Indeed \eqref{rr:eq:model-square} gives the first line of
\eqref{rr:eq:residual-quadratic} with the full Carleman norms.
Their finite orthogonal projections have smaller norm and enter with
a minus sign.  Every other term is exact after compression, since
$J_e$ preserves the step space and both $W$ and $D$ take values there.
Expanding the squared residual proves
\eqref{rr:eq:full-model-residual}.

The intertwining identity now yields
\[
 S_aP_a=I+K_a,\qquad
 S_a=E^*M_a^{-1}E,\qquad K_a=E^*M_a^{-1}\Delta,
 \qquad \|S_a\|\le5.
\]
For every such trial $W,D$, the full correction therefore obeys
\begin{equation}\label{rr:eq:certificate-budget}
 \|K_a\|\le\|E^*W\|_2+
                5\bigl(\mathcal R(W,D)+\|\Delta-D\|_2\bigr).
\end{equation}
This follows by subtracting $E^*W$ and using
$M_a^{-1}\Delta-W=M_a^{-1}(\Delta-D+D-M_aW)$.
The source complement is included in $\Delta-D$.

\subsection{Exact finite quantities}

All finite matrices use normalized cell indicators.  For cells
$(a_i,b_i)$ of length $h_i$, write $F(x)=x\log|x|$, $F(0)=0$.
The Hilbert and Carleman matrices are
\begin{align}
 H_{ij}&=\frac{F(b_i-a_j)-F(a_i-a_j)-F(b_i-b_j)+F(a_i-b_j)}
                   {\pi\sqrt{h_i h_j}},\nonumber\\
 C_{ij}&=\frac{F(b_i+b_j)-F(a_i+b_j)-F(b_i+a_j)+F(a_i+a_j)}
                   {\pi\sqrt{h_i h_j}}.
 \label{rr:eq:double-cell-logarithms}
\end{align}
Since $F''(x)=1/x$ off zero, twice integrating gives these formulas;
the Hilbert diagonal is zero in principal value. The Hilbert formula
also applies to the interior and exterior cells in the existence test.
With the channel structure of $Q_e$ and exchanges $(1,2),(3,4)$,
these matrices evaluate \eqref{rr:eq:residual-quadratic} exactly.

The projected trial norm is equally explicit.  Suppose the source
partition $I_i=(s_i,s_{i+1})$ is reflection symmetric and is also the
initial part of the model partition on $(0,1)$.  Let $v_{j,i}$ be the
constant row vector of $W_j$ on $I_i$, and let $i^*$ index its reflected
cell.  These row vectors are obtained by dividing normalized model-cell
coefficients by $\sqrt{|I_i|}$.  Then
\begin{equation}\label{rr:eq:projected-trial-norm}
 \|E^*W\|_2^2=\sum_i\sum_{j=1}^2
 \left[A_i\|v_{j,i}\|^2+B_i\|v_{j+2,i^*}\|^2
       +2C_i\operatorname{Re}\langle v_{j,i},v_{j+2,i^*}\rangle\right],
\end{equation}
where, with $h_i=s_{i+1}-s_i$,
\[
 A_i=\frac{h_i}{2}+\frac{\sin(\pi s_{i+1})-\sin(\pi s_i)}{2\pi},
 \quad B_i=h_i-A_i,
 \quad C_i=\frac{\cos(\pi s_i)-\cos(\pi s_{i+1})}{2\pi}.
\]
They are the integrals of $\chi^2$, $\chi(1-t)^2$ and
$\chi(t)\chi(1-t)$ on the cell, respectively.

\subsection{A global midpoint error bound}

Take $D$ to be the exact midpoint step kernel of
\eqref{rr:eq:cosine-defect} on $(0,T)\times(0,1)$, zero for $t>T$.
The output partition includes one as an endpoint; the source partition
is exactly reflection symmetric.  For a scalar kernel with derivative
bounds $A,B$ on a rectangle of widths $h,k$, its squared midpoint error
is at most
\begin{equation}\label{rr:eq:midpoint-cell-bound}
 \mathcal E(h,k;A,B)
   =hk\left(\frac{A^2h^2+B^2k^2}{12}+\frac{ABhk}{8}\right).
\end{equation}
To verify this, bound the pointwise error by
$A|t-t_0|+B|s-s_0|$ and integrate its square over the rectangle.

The following bounds specify every cell.  For $t<1$, both derivatives
of $D_0$ are bounded by $c_0^2/2$, by the divided-difference integral
and $|\chi''|\le c_0^2$.  Both derivatives of $D_1$ are bounded by
$3c_0^2/2$: use
$|\chi'(t)|\le c_0^2t$ and
$|\chi(t)-\chi(s)|\le c_0^2|t^2-s^2|/2$ in the differentiated quotient.
For an output cell above one with lower endpoint $t_-$, and source
lower endpoint $s_-$, put $d=t_-+s_-$.  Since
$D_1=-\chi(s)/(t+s)$ there, its derivative bounds are
\[
 |\partial_tD_1|\le d^{-2},\qquad
 |\partial_sD_1|\le c_0d^{-1}+d^{-2}.
\]

The remaining nonzero kernels reduce, up to a sign, to
$f(\alpha,\beta)=\sin(c_0\alpha)/(\alpha+\beta)$,
$0\le\alpha\le1$, $\beta\ge0$.
For $D_2$ below one take $(\alpha,\beta)=(1-t,s)$; for $D_0$ above
one take $(\alpha,\beta)=(1-s,t-1)$.
On $[\alpha_-,\alpha_+]\times[\beta_-,\beta_+]$ with
$d=\alpha_-+\beta_->0$, valid bounds are
\begin{equation}\label{rr:eq:corner-derivative-bounds}
 A_\alpha=\frac{c_0\beta_++c_0^3\alpha_+^3/3}{d^2},
 \qquad A_\beta=\frac{c_0\alpha_+}{d^2}.
\end{equation}
Indeed $|x\cos x-\sin x|\le x^3/3$, by integration of $-x\sin x$.
This proves the first bound after differentiating $f$; the second
uses $|\sin x|\le x$.  Assign these bounds to $t,s$ in the displayed
coordinate order.  When $d=0$, use instead the squared error bound
$c_0^2$ times the cell area, since both $f$ and its midpoint value
lie in $[0,c_0]$.  This handles the corner cells without asserting a
global square-integrable first derivative.  Above one, $D_2$ vanishes.

Let $E_j$ be the sum of these squared cell bounds for $D_j$.
The multiplicities in \eqref{rr:eq:cosine-defect} give the finite error
$(4E_0+2E_1+2E_2)/(4\pi^2)$.
The complete omitted output tail satisfies
\begin{align}
 \|\mathbf1_{(T,\infty)}\Delta\|_2^2
 &=\frac1{4\pi^2}\int_0^1\chi(s)^2
                     \left(\frac4{T-s}+\frac2{T+s}\right)ds\nonumber\\
 &\le\frac1{8\pi^2}\left(\frac4{T-1}+\frac2T\right),
 \label{rr:eq:global-defect-tail}
\end{align}
using $\int_0^1\chi(s)^2ds=1/2$.
Finite and tail squared errors add on disjoint output supports.
These inequalities provide a global bound for $\|\Delta-D\|_2$.

\subsection{Exact witness and rational verification gates}

The source grid has 48 cells and is specified by
\[
 s_j=\frac{j^4}{2\cdot24^4}\quad(0\le j\le24),\qquad
 s_{48-j}=1-s_j.
\]
The model grid is the sorted union of
\[
 \{s_j:0\le j\le48\},\quad
 \{1+s_j:1\le j\le48\},\quad
 \{2^k(1+\ell/6):1\le k\le39,\ 1\le\ell\le6\}.
\]
It has 330 cells and ends at $T=2^{40}$.
The file \texttt{step\_parametrix\_witness.json} specifies a
$1320\times96$ matrix $W$, whose real and imaginary coefficients are
integers divided by $2^{48}$ in the normalized cell bases.
Its SHA-256 digest begins \texttt{2b5d1dca6fbcdf3a}; the full digest is
recorded with the checker output.  The choice of trial is immaterial to
the proof once these exact coefficients are fixed.

The accompanying negative-window checker and its global-error helper
evaluate the preceding expressions with outward-rounded Arb arithmetic
\cite{Johansson2017}; their filenames and execution instructions are
listed in \texttt{anc/README.md}.
For \eqref{rr:eq:midpoint-cell-bound}--\eqref{rr:eq:global-defect-tail}
it uses the rational upper bounds $c_0<11/7$,
$1/(4\pi^2)<1/36$ and $1/(8\pi^2)<1/72$.
Midpoint defect values are enclosed directly, so $D$ has no separate
saved-coefficient rounding error.  All primitive arguments in
\eqref{rr:eq:double-cell-logarithms} are formed as exact rationals
before conversion; a zero argument uses $F(0)=0$.

At both 128 and 192 bits, the checker verifies the strict rational gates
\begin{equation}\label{rr:eq:rational-certificate-gates}
 \|E^*W\|_2<\frac{27}{100},\qquad
 \mathcal R(W,D)<\frac9{500},\qquad
 \|\Delta-D\|_2<\frac{37}{1000}.
\end{equation}
No eigenvalue solve, Fourier truncation, or numerical quadrature enters
this verification.  The gates and \eqref{rr:eq:certificate-budget} give
\[
 \|K_a\|<\frac{27}{100}
          +5\left(\frac9{500}+\frac{37}{1000}\right)
       =\frac{109}{200}<\frac35.
\]

\begin{proof}[Exclusion and uniqueness in Theorem~\ref{boundary:thm:negative-exclusion:rr}]
The Neumann series makes $I+K_a$ invertible.  Since $S_aP_a=I+K_a$
and $\|S_a\|\le5$, the self-adjoint pencil $P_a$ is bounded below
by a constant strictly greater than $2/25$.  Its range is closed and
its orthogonal complement is its zero kernel, so it is onto and
\[
 \|P_a^{-1}\|<\frac{25}{2}.
\]
For real $\lambda$ with $|\lambda-a|\le3/50$,
$\|P_\lambda-P_a\|=|\lambda-a|$ and the inverse Neumann product is
strictly below $3/4$.  Thus every $P_\lambda$ on
$[-953/2000,-713/2000]$ is invertible.  The exact frequency
correspondence of Proposition~\ref{v12:prop:frequency-pencil} excludes
physical spectrum there.

The existence argument above supplies a
positive eigenvalue in $(81/200,107/250)$.  If another positive boundary
level lay in $[713/2000,953/2000]$, the signed alternation of
Corollary~\ref{interaction:cor:boundary-alternation} would put a negative
gap magnitude strictly between them, contradicting the exclusion.
The eigenvalue is simple by Theorem~\ref{boundary:thm:one-cell-flow}.
\end{proof}

